\begin{thebibliography}{99}
\bibitem{ArthurClozel1989}
J. Arthur and L. Clozel, \emph{Simple Algebras, Base Change, and the Advanced Theory of the Trace Formula}, Annals of Mathematics Studies 120, Princeton University Press, 1989. Prime cyclic base change and induction are used with their cuspidality hypotheses.

\bibitem{ClozelRajan2020}
L. Clozel and C. S. Rajan, Solvable base change, \emph{Journal f\"ur die reine und angewandte Mathematik} 772 (2021), 147--174. DOI: \href{https://doi.org/10.1515/crelle-2020-0023}{10.1515/crelle-2020-0023}. Checked preprint \href{https://arxiv.org/abs/1806.02513}{arXiv:1806.02513v1}, 7 June 2018. Theorem 3(2) and the ensuing cyclic descent remark are the fibre inputs.

\bibitem{Rajan2002}
C. S. Rajan, On the image and fibres of solvable base change, \emph{Mathematical Research Letters} 9 (2002), 499--508.

\bibitem{KimYamauchi2016}
H. H. Kim and T. Yamauchi, A uniform structure on subgroups of \(GL_n(\mathbb F_q)\) and its application to a conditional construction of Artin representations of \(GL_n\), \emph{Journal of the Ramanujan Mathematical Society} 32(1) (2017), 75--99. Checked preprint \href{https://arxiv.org/abs/1403.5535}{arXiv:1403.5535v3}, 7 April 2016. Theorem 1.1 is conditional on its five inputs. The inherited citation key records the preprint year.

\bibitem{Ramakrishnan2002}
D. Ramakrishnan, Modularity of solvable Artin representations of \(GO(4)\) type, \emph{International Mathematics Research Notices} 2002(1), 1--54. DOI: \href{https://doi.org/10.1155/S1073792802000016}{10.1155/S1073792802000016}.

\bibitem{Langlands1980}
R. P. Langlands, \emph{Base Change for \(GL(2)\)}, Annals of Mathematics Studies 96, Princeton University Press, 1980.

\bibitem{Tunnell1981}
J. Tunnell, Artin's conjecture for representations of octahedral type, \emph{Bulletin of the American Mathematical Society} 5(2) (1981), 173--175.

\bibitem{Neukirch1999}
J. Neukirch, \emph{Algebraic Number Theory}, Springer, 1999.

\bibitem{Henniart2000}
G. Henniart, Une preuve simple des conjectures de Langlands pour \(GL(n)\) sur un corps p-adique, \emph{Inventiones Mathematicae} 139(2) (2000), 439--455.

\bibitem{HarrisTaylor2001}
M. Harris and R. Taylor, \emph{The Geometry and Cohomology of Some Simple Shimura Varieties}, Annals of Mathematics Studies 151, Princeton University Press, 2001. Checked author preprint dated 30 August 1999, Lemma 12.7(4--6), for twisting, cyclic prime base change, and induction compatibility. Its numbering differs from the published book.

\bibitem{Brauer1947}
R. Brauer, On Artin's \(L\)-series with general group characters, \emph{Annals of Mathematics} 48 (1947), 502--514.

\bibitem{Konig2010}
J. K\"onig, Solvability of generalized monomial groups, \emph{Journal of Group Theory} 13(2) (2010).

\bibitem{CimpoeasNicolae2020}
M. Cimpoea\c{s} and F. Nicolae, Artin \(L\)-functions of almost monomial Galois groups, \emph{Forum Mathematicum} 32(4) (2020).

\bibitem{Brown1982}
K. S. Brown, \emph{Cohomology of Groups}, Graduate Texts in Mathematics 87, Springer, 1982.

\bibitem{Isaacs1976}
I. M. Isaacs, \emph{Character Theory of Finite Groups}, Academic Press, 1976.

\bibitem{Chenevier2014}
G. Chenevier, The p-adic analytic space of pseudocharacters of a profinite group and pseudorepresentations over arbitrary rings, in \emph{Automorphic Forms and Galois Representations}, Vol. 1, London Mathematical Society Lecture Note Series 414, Cambridge University Press, 2014, 221--285. DOI: \href{https://doi.org/10.1017/CBO9781107446335.008}{10.1017/CBO9781107446335.008}. Checked \href{https://arxiv.org/abs/0809.0415}{arXiv:0809.0415v2}, 18 July 2013: Proposition 1.27, Theorem 2.12, Corollary 2.13, and Section 4.1. These are preprint numbering. The algebraic realization theorem alone does not assert continuity for the usual complex or l-adic topology.

\bibitem{Scholze2015}
P. Scholze, On torsion in the cohomology of locally symmetric varieties, \emph{Annals of Mathematics} 182(3) (2015), 945--1066. DOI: \href{https://doi.org/10.4007/annals.2015.182.3.3}{10.4007/annals.2015.182.3.3}. Corollary 5.4.3, p. 1060, is the residual eigensystem input. Field conditions, cohomological occurrence, normalization, and the source's historical dependencies are retained.

\bibitem{GoldringShin2014}
W. Goldring, with an appendix by S. W. Shin, Galois representations associated to holomorphic limits of discrete series, \emph{Compositio Mathematica} 150(2) (2014), 191--228. DOI: \href{https://doi.org/10.1112/S0010437X13007355}{10.1112/S0010437X13007355}. Theorem 1.2.1 is used only with its unitary datum, coefficient prime, and split prime normalization conditions.

\bibitem{GeeKassaei2013}
T. Gee and P. Kassaei, Companion forms in parallel weight one, \emph{Compositio Mathematica} 149(6) (2013), 903--913. DOI: \href{https://doi.org/10.1112/S0010437X12000875}{10.1112/S0010437X12000875}. Checked \href{https://arxiv.org/abs/1206.6631}{arXiv:1206.6631v1}, 28 June 2012. Its companion and Artin consequences retain their respective modularity hypotheses.

\bibitem{DimitrovWiese2020}
M. Dimitrov and G. Wiese, Unramifiedness of Galois representations attached to Hilbert modular forms mod \(p\) of weight one, \emph{Journal of the Institute of Mathematics of Jussieu} 19(2) (2020), 281--306. DOI: \href{https://doi.org/10.1017/S1474748018000026}{10.1017/S1474748018000026}. Checked preprint version 4, 28 November 2017. The full Frobenius polynomial comparison also requires the diamond character determinant normalization.

\bibitem{BrauerNesbitt1941}
R. Brauer and C. Nesbitt, On the modular characters of groups, \emph{Annals of Mathematics} 42(2) (1941), 556--590.

\bibitem{CurtisReiner1962}
C. W. Curtis and I. Reiner, \emph{Representation Theory of Finite Groups and Associative Algebras}, John Wiley and Sons, 1962.

\bibitem{Serre1977}
J.-P. Serre, \emph{Linear Representations of Finite Groups}, Graduate Texts in Mathematics 42, Springer, 1977.

\bibitem{Deligne1973}
P. Deligne, Les constantes des \'equations fonctionnelles des fonctions \(L\), in \emph{Modular Functions of One Variable II}, Lecture Notes in Mathematics 349, Springer, 1973, 501--597.

\bibitem{DeligneHenniart1981}
P. Deligne and G. Henniart, Sur la variation, par torsion, des constantes locales d'\'equations fonctionnelles de fonctions \(L\), \emph{Inventiones Mathematicae} 64(1) (1981), 89--118. Theorem 4.6, p. 110, supplies the stable epsilon ratio used here.

\bibitem{BLGGT2014}
T. Barnet-Lamb, T. Gee, D. Geraghty, and R. Taylor, Potential automorphy and change of weight, \emph{Annals of Mathematics} 179(2) (2014), 501--609. DOI: \href{https://doi.org/10.4007/annals.2014.179.2.3}{10.4007/annals.2014.179.2.3}. Checked \href{https://arxiv.org/abs/1010.2561}{arXiv:1010.2561v4}, 9 December 2013. Theorem 4.2.1 uses regular algebraic and further lifting hypotheses.

\bibitem{CalegariGeraghty2020}
F. Calegari and D. Geraghty, Minimal modularity lifting for nonregular symplectic representations, \emph{Duke Mathematical Journal} 169(5) (2020), 801--896. DOI: \href{https://doi.org/10.1215/00127094-2019-0044}{10.1215/00127094-2019-0044}. This is a special nonregular range, not the arbitrary all zero weight gate.

\bibitem{DeligneSerre1974}
P. Deligne and J.-P. Serre, Formes modulaires de poids 1, \emph{Annales scientifiques de l'\'Ecole Normale Sup\'erieure} 7(4) (1974), 507--530. DOI: \href{https://doi.org/10.24033/asens.1277}{10.24033/asens.1277}. Lemma 3.2, p. 513, is a primary instance of the Frobenius determination argument.

\bibitem{Wang2026}
J. Wang, Strong Artin conjecture for generalized octahedral representations in \(GL_3\), \href{https://arxiv.org/abs/2609.38231}{arXiv:2609.38231v1}, 28 September 2026. Theorem 1.1 and Corollary 5.4 state the special three and six dimensional cases mentioned here. Statements checked. The full proof has not been independently audited for this reader edition.
\end{thebibliography}
